\documentclass{tufte-handout}

\title{The Map That Ate the Territory}
\author{J. Okonkwo-Reyes}
\date{}

\usepackage{booktabs}

\begin{document}

\maketitle

\begin{abstract}
Metrics dashboards are supposed to represent an organisation's work.
This essay argues that past a certain size, the dashboard stops
representing the work and starts \emph{replacing} it as the object of
attention --- and that the failure mode is structural, not a matter
of picking better metrics.
\end{abstract}

Every growing organisation eventually builds a dashboard: a single
screen distilling the state of the business into a handful of
numbers a leader can absorb in ten seconds.\sidenote{The instinct is
sound. A CEO cannot personally read every support ticket, every
commit, every customer call --- a dashboard is a necessary
compression of information that would otherwise be inaccessible at
scale.} The trouble begins once the dashboard stops being a
\emph{window} onto the work and becomes the thing people manage
\emph{to}.

\section{The compression is lossy by construction}

A dashboard number is, definitionally, a lossy compression of
something richer. ``Support tickets resolved per day'' compresses
away whether the resolution actually helped the customer. ``Commits
per engineer per week'' compresses away whether any given commit made
the codebase better or worse. This is not a flaw that better metric
design can fix --- it is what a metric \emph{is}: a projection from a
high-dimensional reality onto one number, and every projection
discards information along whichever dimensions it does not
capture.\sidenote{This is the same shape of argument mathematicians
make about projections in linear algebra: a projection onto a
subspace necessarily loses whatever lay outside that subspace, no
matter how well-chosen the subspace is.}

\begin{marginfigure}
\begin{tabular}{ll}
\toprule
\textbf{Reality} & \textbf{Dashboard} \\
\midrule
customer helped & ticket closed \\
code improved & commit merged \\
user retained & session logged \\
\bottomrule
\end{tabular}
\caption{Each dashboard column discards everything in its row except
the one thing it counts.}
\end{marginfigure}

\section{Optimisation pressure finds the seam}

The lossiness would be harmless if nobody optimised against the
compressed number specifically. But once a team's performance is
\emph{evaluated} on ``tickets resolved per day'', the fastest way to
raise that number is not necessarily to help customers faster --- it
is to find whatever behaviour raises the count with the least effort,
which is exactly the behaviour the compression was blind to in the
first place. The dashboard did not merely fail to capture the
seam between the metric and the goal; its existence as a target
actively widens that seam, because now there is a reason to search
for it.

\section{Why "add more metrics" does not fix this}

The standard response to a gamed metric is to add a second metric
that catches the gaming --- customer satisfaction alongside ticket
count, code review approval alongside commit count. This narrows the
seam but cannot close it, for the same reason one projection plus a
second projection still is not the full-dimensional reality: two
numbers are still two numbers, and an organisation with enough scale
to need a dashboard in the first place has a process too
high-dimensional for any small fixed set of numbers to fully pin
down.\sidenote{This is why the healthiest organisations tend to pair
dashboards with judgment --- a manager who actually reads a sample of
the underlying tickets, not just the count --- rather than trying to
solve the problem with a bigger dashboard.}

\section{The map is not obligated to become the territory}

None of this is an argument against dashboards. It is an argument
against treating the dashboard as a sufficient substitute for
looking at the underlying work. The map is useful precisely because
it is smaller than the territory; the failure mode is forgetting
that a smaller thing was ever a compression at all, and starting to
manage the map as though it were the territory itself.

\end{document}
